\noindent\textbf{Resultado.} Con \(t\) en minutos y cantidades en kg,
\[A(t)=3e^{-0.04t}+2e^{-0.12t},\qquad B(t)=6e^{-0.04t}-4e^{-0.12t}.\]

\noindent\textit{Desarrollo.} Las variables de estado son las masas \(A(t),B(t)\) de sal, en kg. Los balances
de volumen son \(6+2-8=0\) y \(8-2-6=0\), por lo que ambos volúmenes permanecen en 100 L.
La mezcla uniforme permite usar concentraciones \(A/100\) y \(B/100\), en kg/L.
El balance entrada menos salida proporciona
\[A'=-0.08A+0.02B,\qquad B'=0.08A-0.08B,\qquad A(0)=5,\ B(0)=2.\]
Los coeficientes tienen unidades min\(^{-1}\). Antes de resolver,
\((A+B)'=-0.06B\): solo la salida externa elimina sal.

Los autovalores son \(-0.04\) y \(-0.12\), con vectores \((1,2)^T\) y \((1,-2)^T\).
Así, \((A,B)^T=c_1e^{-0.04t}(1,2)^T+c_2e^{-0.12t}(1,-2)^T\).
Las condiciones iniciales dan \(c_1=3,c_2=2\).
La sustitución verifica las dos ecuaciones y las condiciones.
Para \(t\ge0\), \(A>0\) y \(B=2e^{-0.04t}(3-2e^{-0.08t})>0\).
La sal total decrece hacia cero, aunque B crece inicialmente porque \(B'(0)=0.24\) kg/min.
Ambas masas tienden a cero, como corresponde a la renovación con agua pura.
El resultado presupone caudales constantes, ausencia de pérdidas de sal distintas de las descritas
y mezcla instantánea; no representa gradientes locales ni tiempos de agitación.
